\begin{statement}{A.1}{Open}{The two principal assertions}
The classification of all noncuspidal continuous sections of $E_U\to G_\Q$
by $U(\Q)$, and the lifting of every such section along the continuous
surjection $G_{\Q(t)}\to E_U$, remain unresolved in this manuscript.
\end{statement}
These are precisely items A1(1) and A1(2), corresponding respectively to
cote 149, scan 56, and scan 127, conjecture b). We do not treat a central
lift, a finite quotient, or a Kummer class as a proof of either assertion.

\begin{statement}{A.2}{Proved}{Local multiplicative completions}
For every finite prime $p$, the canonical homomorphism
$\Q_p^\times\to\widehat{\Q_p^\times}$ is injective.
The canonical homomorphism $\R^\times\to\widehat{\R^\times}$
has kernel $\R_{>0}$ and identifies its target with $\{\pm1\}$.
\end{statement}
\begin{proof}
Write $a=p^m u$ with $m\in\Z$ and $u\in\Z_p^\times$.
An element lying in every power subgroup has $m\in\bigcap_n n\Z=0$.
For odd $p$, the decomposition
$\Z_p^\times=\mu_{p-1}\times(1+p\Z_p)$ and the $p$-adic logarithm
identify the second factor with $p\Z_p$. Taking powers divisible by
$(p-1)p^k$ annihilates the first factor and makes the logarithm divisible
by arbitrarily large powers of $p$. Thus $u=1$.
For $p=2$, use $\Z_2^\times=\{\pm1\}\times(1+4\Z_2)$ and the
logarithm on the second factor. Even powers remove the sign and powers
of 2 again give trivial intersection. This proves injectivity at every
finite prime, including 2. Positive real numbers have roots of every
positive integral order, while even powers of nonzero real numbers are
positive. The compatible quotients therefore retain exactly the sign.
\end{proof}

\begin{statement}{A.3}{Refuted}{Statement (G) as printed}
The assertion (G) of L2 and A12 is false, even with every $u_c$ represented
by a nonzero rational number.
\end{statement}
\begin{proof}
Define $x_2=\infty$ and $x_v=2$ for every other finite or real place $v$.
For $c\in\Q\setminus\{2\}$, let $u_c$ be the image of $2-c$ in
$\widehat{\Q^\times}$, and let $u_2=1$.
For $c\ne2$, every required equality is the image of the equality
$x_v-c=2-c$ at a place $v\ne2$. At $v=2$ the condition in (G) omits
the comparison because $x_2=\infty$. For $c=2$ it omits every place,
because each coordinate is either $c$ or $\infty$.
The tuple is not the diagonal tuple of any element of $\Pj^1(\Q)$:
its coordinates at 2 and 3 are different. All maps used here are the
homomorphisms of exponent completions induced by the field embeddings
$\Q\hookrightarrow\Q_v$.
\end{proof}
This counterexample is a tuple of local points and Kummer classes. It
does not construct a continuous section of $E_U\to G_\Q$.

\begin{statement}{A.4}{Refuted}{Every finite puncture test in A12}
For each finite $S\subset\Q$, the version of (G) testing only $c\in S$
admits a nondiagonal tuple with $x_p\in U(\Q_p)$ at every finite prime
and all $u_c$ rational. Consequently no finite $S\supset\{0,1\}$
has the property requested in A12.
\end{statement}
\begin{proof}
Choose $r\in\Q\setminus(S\cup\{0,1\})$. The complement
$\R\setminus S$ has an open interval containing $r$. Choose an
irrational $y$ in that interval and put $x_p=r$ for every finite prime,
$x_\infty=y$, and $u_c=r-c$ for $c\in S$.
At finite primes the equality is literal. At the real place $y-c$
and $r-c$ have the same sign for every $c\in S$, so A.2 gives equality
of their completed Kummer classes. The real coordinate prevents the
tuple from being a rational diagonal. No finite coordinate is a cusp.
\end{proof}

\begin{statement}{A.5}{Proved}{The exact content of one rational Kummer class}
Fix $c\in\Q$, $r\in\Q^\times$, and a tuple $(x_v)\in\prod_v\Pj^1(\Q_v)$.
The comparison demanded in (G) for $u_c=r$ holds exactly when
$x_p\in\{c,c+r,\infty\}$ at each finite prime $p$, and at the real
place either $x_\infty\in\{c,\infty\}$ or
$\operatorname{sgn}(x_\infty-c)=\operatorname{sgn}(r)$.
\end{statement}
\begin{proof}
At a finite place where the comparison is required, A.2 turns equality
in the completion into $x_p-c=r$. At the omitted places the displayed
three-point condition simply records the omission. At the real place
A.2 turns equality into equality of signs. These steps are reversible.
\end{proof}

\begin{statement}{A.6}{Proved}{The entire rational-class part of (G)}
For a tuple satisfying (G) with every $u_c$ represented by a nonzero
rational number, either all finite coordinates are $\infty$, or there
is a unique $r\in\Q$ such that every finite coordinate is in
$\{r,\infty\}$. In the second case the real coordinate is also in
$\{r,\infty\}$. Conversely every such mixed tuple satisfies (G);
when all finite coordinates are $\infty$, any real coordinate is
permitted with suitable rational choices of $u_c$.
\end{statement}
\begin{proof}
Suppose $x_p$ is finite at one prime. Choose a rational $c\ne x_p$.
A.2 gives $x_p=c+u_c\in\Q$; call it $r$. For every rational
$c\ne r$, injectivity at that prime forces $u_c=r-c$.
If another finite coordinate is a finite number $z\ne r$, choose
a rational $c$ different from both $z,r$; the same equality gives
$z=r$, a contradiction. Thus all finite coordinates are $r$ or
$\infty$. For a finite real coordinate $y\ne r$, choose a rational
$c$ strictly between $y$ and $r$. The signs of $y-c$ and $r-c$
are opposite, contradicting (G). This proves the assertion, including
uniqueness when a finite coordinate is present.
Conversely set $u_c=r-c$ for $c\ne r$ and $u_r=1$.
Every required equality holds. If all finite coordinates are
$\infty$, only a real sign condition remains; choose $u_c=1$ or
$-1$ according to that sign, using any value at an omitted coordinate.
\end{proof}

\subsection{The actual maps in the self-cover leads}

\begin{statement}{A.7}{Proved}{All nonconstant self-morphisms of the tripod}
Over a characteristic-zero field, every nonconstant morphism
$U\to U$ is one of the six automorphisms permuting
$\{0,1,\infty\}$. In particular its degree is one.
\end{statement}
\begin{proof}
Base change to an algebraic closure and extend the morphism uniquely
to a rational map $f:\Pj^1\to\Pj^1$ of degree $d$.
Every inverse image of $0,1,\infty$ must belong to the three-point
boundary, since the given restriction lands in $U$.
Each of the three fibres has degree $d$, counting ramification.
The ramification contribution of these fibres is therefore at least
$3d-3$. Riemann--Hurwitz gives total ramification $2d-2$,
so $3d-3\leq2d-2$ and $d\leq1$. A degree-one map with the stated
boundary property permutes those three points and is one of
\[
 t,\quad 1-t,\quad t^{-1},\quad (1-t)^{-1},
 \quad t/(t-1),\quad (t-1)/t .
\]
All six are defined over the prime field. This lists the maps as
morphisms, not just their induced outer homomorphisms.
\end{proof}

\begin{statement}{A.8}{Refuted}{The proposed higher-degree self-cover action}
The maps $4t(1-t)$ and $j(t)/1728$ do not define self-morphisms of
$U$. Their actual finite \'etale restrictions have degrees two
and six and domains
\[
 V_2=\Pj^1\setminus\{0,1,\tfrac12,\infty\},
 \qquad
 V_6=U\setminus\{-1,\tfrac12,2,\zeta_6,\zeta_6^{-1}\},
\]
respectively. Consequently L12(iv) and L13 do not define the
asserted semigroup action on $\Sec(U/\Q)$ by these maps.
\end{statement}
\begin{proof}
For $f=4t(1-t)$, the fibre of 1 is $2[\tfrac12]$, while the
fibres of 0 and $\infty$ have supports $\{0,1\}$ and $\{\infty\}$.
Its only critical points are $\tfrac12,\infty$.
For $J=j/1728$ we have
\[
 J=\frac{4(t^2-t+1)^3}{27t^2(t-1)^2},\qquad
 J-1=\frac{(t-2)^2(t+1)^2(2t-1)^2}{27t^2(t-1)^2}.
\]
The derivative has numerator proportional to
$(t-2)(t+1)(2t-1)(t^2-t+1)^2$.
These formulae identify precisely the inverse images of the removed
target points and show that all ramification lies over them.
The restrictions stated above are therefore finite \'etale.
They give the continuous homomorphisms
\[
 \pi_1(V_i)\xrightarrow{f_{i*}}\pi_1(U),
 \qquad
 \pi_1(V_i)\xrightarrow{j_{i*}}\pi_1(U),
\]
where the first is an open injection and the second comes from the
open immersion and is surjective. A section of $\pi_1(U)$ supplies
no inverse for the second homomorphism merely by functoriality.
A chosen full lift to $\pi_1(V_i)$ can be composed with $f_{i*}$;
this describes the actual correspondence on sections.
A.7 also rules out any further higher-degree rational self-morphism
of $U$. Thus infinitely many distinct such self-morphisms, as used
in the stated fixed-section test, do not exist.
\end{proof}
The formulae and derivative factorizations are checked exactly in
\texttt{code/root\_checks.py}. An action through additional lifts
or a specified orbifold would require its own domain and comparison
maps; no claim about that unformulated action is made here.

\subsection{The coefficient prime in the local cusp test}

\begin{statement}{A.9}{Refuted}{L12(i) with every coefficient prime}
The ordinary rational point section at $t=5$ gives, at the
ramified prime 5, a local geometric pro-5 section outside every
cuspidal packet. This already holds in its mod-5 geometric
abelianization.
\end{statement}
\begin{proof}
The maps $t,1-t:U_{\Q_5}\to\G_m$ induce the two Kummer coordinates
in $H^1(G_{\Q_5},\mu_5)^2$. At the point 5 they are
\[
                 \bigl(\kappa_5(5),\kappa_5(-4)\bigr).
\]
Their valuation coordinates modulo 5 are 1 and 0. The fifth powers
of units modulo 25 are exactly $\{1,7,18,24\}$, whereas
$-4\equiv21\pmod {25}$. Thus the second Kummer class is nonzero.
At cusp 0 the second class is zero; at cusp 1 the first is zero;
at infinity the two classes are equal because the ratio
$(1-t)/t$ specializes to $-1$, itself a fifth power.
The displayed pair satisfies none of these necessary conditions,
so it cannot be the image of a cuspidal pro-5 section.

This is a place of genuine geometric ramification: for $E_5$,
$\Delta=6400$ has 5-adic valuation 2 and $c_4=336$ is a unit.
The reduction is a nodal cubic with tangent slopes
$\pm\sqrt{-1}$; these lie in $\Q_5$. Thus it is split
multiplicative, with nonzero tame monodromy in every Tate module
of coefficient prime different from 5.
The finite fifth-power certificate is enumerated in the code.
\end{proof}
At a full locally cuspidal section, restriction of the inclusion of
its cusp decomposition group to each pro-$\ell$ quotient does give
compatible images of its one full Kummer parameter. That
functorial assertion is expressed by the maps
$\widehat{\Z}(1)\to\Z_\ell(1)$. A.9 proves that local cuspidality
in every pro-$\ell$ quotient at a ramification place does not
hold for all actual global sections.

\subsection{The completed multiplicative group and the unit equation}

\begin{statement}{A.10}{Proved}{Explicit valuation coordinates}
There is an isomorphism of topological abelian groups
\[
 \widehat{\Q^\times}\simeq
 \{\pm1\}\times
 \left\{(a_q)_q\in\prod_q\widehat{\Z}:
    \text{for every }n,\ a_q\in n\widehat{\Z}
                  \text{ for all but finitely many }q\right\}.
\]
It is induced by the sign and the prime valuations.
The group on the right is given the inverse-limit topology: a basic
neighbourhood of zero requires every coordinate, simultaneously, to
vanish modulo a fixed $n$. This specifies the topology used in the
isomorphism.
A class is locally represented by an element of $\Q_p^\times$
exactly when its $p$-valuation belongs to $\Z\subset\widehat{\Z}$.
It is represented by a rational number exactly when all these
valuations are integers and only finitely many are nonzero.
\end{statement}
\begin{proof}
Unique factorization gives
$\Q^\times=\{\pm1\}\times\bigoplus_q q^{\Z}$.
Taking the quotient by $n$th powers and then the inverse limit
gives precisely the displayed eventual-divisibility condition:
at each modulus the valuation vector has finite support.
Conversely such vectors define compatible finite-support classes
at every modulus, establishing surjectivity and the inverse map.
Both maps are continuous by their finite-level descriptions.
For comparison with the product topology on
$\prod_q\widehat{\Z}$, the identity map from this inverse-limit
topology to the subspace product topology is continuous. Its inverse
is not continuous: vectors with a single coordinate equal to 1 at
successively different primes converge to zero coordinatewise, while
their reductions modulo 2 never equal the zero vector. The
simultaneous finite-support quotients therefore must be retained.

The local decomposition
$\Q_p^\times=p^\Z\times\Z_p^\times$ gives
$\widehat{\Q_p^\times}=p^{\widehat{\Z}}\times\Z_p^\times$,
since the compact unit group is complete for its power quotients.
Thus the local image is an actual nonzero number precisely for an
integer valuation. The final assertion follows from the global
isomorphism: a finitely supported integer vector and a sign specify
one rational number, and unique factorization proves the converse.
\end{proof}

\begin{statement}{A.11}{Proved}{Infinite-support ordinary solutions exist}
The ordinary local completed unit equation has pairs whose two
valuation vectors have infinite support. An explicit recursive
construction, with actual nonzero $p$-adic local values at every finite
prime and cofinally divisible ordinary integer valuations, is given
and proved in A.44. Thus the existence assertion in L9 is true.
\end{statement}
\begin{proof}
Apply A.44 to the two sequences of primes constructed there. Its
valuation table identifies both infinite supports, and its local
limit argument proves the ordinary unit equation at every finite
prime. These are exactly the data in the definition of $\mathcal G$.
\end{proof}
A.3 and A.4 do not settle this assertion: their Kummer classes have
finite support, and A.3 also uses an omitted coordinate.
For an arbitrary pair in A.10, requiring the local equality
$x_p+(1-x_p)=1$ is an additional equation at every prime.
An independent choice of two eventually divisible valuation
vectors does not establish those equations.

At modulus $n$, choose representatives $a_n,b_n\in\Q^\times$.
The equations demanded by an actual solution are exactly local
solubility of
\[
             a_n X^n+b_n Y^n=1,\qquad XY\ne0,
\]
with compatible local roots as $n$ varies. Transition to a larger
modulus requires a twist above the previously chosen twist with
local points at every place. Weak approximation supplies finitely
many prescribed local conditions. The prime-selection construction in
A.44 supplies the additional conditions at every newly introduced
coefficient prime and a compatible infinite tower with infinite support.
A.48 gives the corresponding construction for every fixed finite
rational puncture set. The ordinary all-rational-punctures extension
requires an additional initialization step, recorded after A.48.

The full geometric abelianization
$\Pi\to\Pi^{\rm ab}\simeq\widehat{\Z}(1)^2$ sends a section to
$(A_s,B_s)$ by pullback of $t,1-t$. The actual class-two quotient
and its central embedding problem are identified in A.38, where
A.27 gives the vanishing of the full continuous cup for every
ordinary local pair. A.43 and A.46 compute the next obstruction,
including its derived-limit component. A.52 proves survival through
every prescribed finite third layer for a suitably constructed ghost.
A.58--A.59 use a further prime recursion to construct a single
infinitely supported ordinary pair with an actual full integral
class-three lift. Its passage to the metabelian quotient and to
the full fundamental group remains open.
